\section{Reconciliation with the Budget Lab}\label{sec:reconciliation}

The two models measure different totals: theirs is federal individual income
tax plus payroll plus a corporate wedge, with no states and no non-credit
transfers; this paper's household total spans federal and state and nets out
refundable credits and the benefit programs counted in household net income,
a set that excludes health programs, Head Start and Early Head Start
(Section~\ref{sec:method}). Setting the headline numbers side by side without
a bridge would be wrong, so this section builds one, using only this paper's
run output and objects they publish. Their repository commits the complete
published result grid --- revenue by instrument for all eighteen cells ---
which makes the comparison possible at the cell level rather than against
two text-stated totals.

\textbf{Basis.} Their microsimulation total is individual income tax plus
payroll minus tax-credit outlays; the matching quantity here is federal
income tax net of refundable credits, plus payroll (both halves and
self-employment). Their corporate wedge reduces to the capital growth rate
times CBO's 2030 corporate anchor of \$486B --- in their formula
$\Delta R_{\mathrm{CIT}} = X \cdot (R_{\mathrm{CIT}}/Y^0_K)$ with
$X = Y^0_K\,g$, the capital base cancels --- and all \genTheirNCells{}
committed values equal $g \times \$486\mathrm{B}$ to within
\$\genCitWedgeMaxDevB B, checked against their own per-cell capital growth
rates. Adding their wedge
to this paper's federal microsimulation total yields a like-for-like
``bridged'' headline. The wedge is not PolicyEngine modelling corporate tax,
and it inherits every constant-rate assumption they list.

\begin{table}[t]
\centering
\tabBridge
\caption{The bridged headline cell (Rapid~/~expansive, their central
published scenario), \$B, 2030. ``IIT (net)'' is individual income tax net
of refundable credits; the corporate wedge is theirs in both rows.}
\label{tab:bridge}
\end{table}

\textbf{Where the models agree.} Payroll responses agree closely in all nine
as-forecast cells --- the largest gap anywhere in the grid is
\$\genPayrollMaxCellGap B (Table~\ref{tab:cells}). Across the Rapid variants
ours run \genPayrollTripletOurs{} against theirs
\genPayrollTripletTheirs{} for compressive~/~proportional~/~expansive ---
two independent models putting similar wage distributions against the same
taxable maximum, including the sign reversal under compression. The tilt
ratios also line up once stated on the same basis: on their
federal-plus-corporate accounting, shares-fixed revenue exceeds as-forecast
revenue by
\genTiltOursMatchedSlow$\times$~/~\genTiltOursMatchedMod$\times$~/~\genTiltOursMatchedRapid$\times$
here against
\genTiltTheirsSlow$\times$~/~\genTiltTheirsMod$\times$~/~\genTiltTheirsRapid$\times$
in their grid for Slow~/~Moderate~/~Rapid --- agreement on the paper's
central mechanism, with the same ordering. (The household-total ratio for
Rapid is larger, \genTiltOursHouseholdRapid$\times$, mainly because the
household basis omits the corporate wedge: the wedge grows with the capital
shock, so it cushions the as-forecast denominator on the matched basis.
States and transfers raise the tilt's dollar cost but leave the ratio
essentially unchanged.) Their most extreme published comparison
reproduces in direction but not in magnitude: for Slow with compressive
inequality, their grid puts revenue \genSlowCompTheirsLowerPct\% below the
shares-and-distribution-unchanged counterfactual (their text states 82\%),
while this paper gets \genSlowCompOursLowerPct\%. Slow cells are
also the least stable in this model: Section~\ref{sec:stability} shows them
moving up to \genStabMaxRelPct\% on a data-build refresh alone, because
their bases are single-digit billions. The agreement to take seriously in
this scenario is the sign and the ordering, not the second digit.

\begin{table}[t]
\centering
\small
\tabPayrollCells
\caption{Cell-level comparison on the matched federal basis, all nine
as-forecast cells, \$B, 2030.}
\label{tab:cells}
\end{table}

\textbf{Where they diverge, and why.} The bridged headline is
\genBridgeOursTotal B against their \genBridgeTheirsTotal B ---
\genBridgeRatio$\times$, or \genBridgeOursCboPct\% against
\genBridgeTheirsCboPct\% of the CBO 2030 revenue baseline
(\$\genTheirCBORevB B, from their committed grid). The gap concentrates in
the income-tax response to the capital shock, and the absolute gap says so
more clearly than any ratio: it is nearly constant at
\genIITGapTriplet{}~\$B across the three Rapid wage-spread cells, exactly
what a divergence in the shared capital shock --- rather than in the varying
wage channel --- would produce. (In ratio terms:
\genIITRatioProp$\times$ theirs on the proportional cell, where the
wage-spread channel is neutral; \genIITRatio$\times$ on the expansive
bridge cell.) Three model differences point the same direction. First, the
realized-capital base: their committed cell parameters put their baseline
realized taxable capital base at \$\genTheirYZeroKT T, against
\$\genModelledCapitalShareT T here --- \genBaseGapPct\% larger. Second,
allocation: they distribute the incremental capital flow by SCF-imputed
total wealth; this paper distributes by existing realized capital income,
which they note is the more top-concentrated choice --- and
top-concentrated dollars face higher marginal rates. Third, retirement
routing: taxable pension and IRA distributions here scale with the capital
shock and are taxed at ordinary rates. The capital-scope sensitivity of
Section~\ref{sec:scope} shows that last channel alone is worth about
\$\genRapidScopeDiffAbs B in Rapid.
